\scanpage{3}
\noindent\hfill\fbox{\unsure{Shimura} - Koizumi}

\uline{Proposition} 1 Soient $V$ un anneau de valuation
discrète, \inserted{de corps $K$} $X$ un \unsure{préschém.|préschéma} sur $V$, $F$ un faisceau
quasi-cohérent sur $X$, \inserted{$X'$ l'\unsure{ouvert} $\unsure{X_K|X}$ de $X$, $F'=F_K=F\vert{}X'$}
\deleted{\unsure{Pour} \unsure{que|qu.} \unsure{quotient} \unsure{Pour} que $F$ soit plat sur
$V$, il f. et s. que le plus petit \unsure{de schém.|sous-schéma}
\unsure{quo.|q.} coh. $G$ de $F$ soit plat sur $V$, il \unsure{f}}

\deleted{\unsure{À t}} (i) Pour tout faisceau quotient q. coh.
$G'=F'_{\deleted{\unsure{K}}}/J'$, soit $G$ le faisceau
quotient \deleted{de} $F/J^*$, où $J^*$ est le plus
\deleted{\unsure{grand|plus} petit \unsure{faisceau}}\inserted{grand ss-Module} de $F$ induisant \deleted{$J'$}
\deleted{sur} sur $X'=X_K$ un \deleted{\unsure{sous}} Module $\subset J'$, --
donc $G$ est le plus petit \inserted{Module} quotient de $F$
induisant sur \unsure{$\vert$} $X'_{\deleted{\unsure{K}}}$ un \deleted{\unsure{sous}} faisceau
quotient de $F'$ \deleted{plus grand} majorant $G'$.
Alors $G$ est quasi-cohérent et $V$-plat, et
\[
G|X'=G'
\]

(ii) \deleted{Soit $G=F/$ \unsure{Le} \unsure{un}} Soit $G=F/J$ un
faisceau q. coh de \deleted{\unsure{$G$|$F$}}$F$ qui est plat sur $V$,
et soit $G'=G|X'$ le \unsure{qu.|quotient}\deleted{\unsure{t}} correspondant
de $F'$. Alors $G$ est le plus petit\deleted{\unsure{e}}
\deleted{\unsure{M.}}\unsure{ent} quotient de \deleted{\unsure{$G$}}$F$ \deleted{\unsure{majorant}} qui
induise $G'$ sur $X'$.

\uline{Corollaire 1} Correspondance biunivoque entre
l'ensemble des faisceaux q. coh. \deleted{\unsure{plats}} de
$F$ qui sont $V$-plats, et des faisceaux \inserted{quotients de $F'$} q. coh.
% archive sheet number: 2

